Interpretability claims are often evaluated by accumulation: more correlations, more visual
examples, more downstream benchmarks.  A typed framework suggests a different standard.  Every
semantic upgrade should have an observation that could reject it while leaving weaker claims
intact.

The first level is mathematical calibration.  A gauge projection should satisfy its zero-marginal
identities.  Reference transport should preserve the underlying equivalence class.  A Potts audit
should recover the known coupling in the additive limit.  A claimed Hopfield realization should
show Lyapunov descent under the conditions of the proof.  These tests establish whether the named
mathematical object has actually been implemented; biological validation begins at a later rung.

The second level compares estimators.  Static MSA fields and mean pLM responses should be evaluated
as full categorical blocks in a common reference, with their estimand gap and uncertainty exposed.
Dynamic residuals should be tested across natural, mutational, and held-out backgrounds.  Route
reversibility should be separated from response reciprocity.  Null ensembles should preserve the
nuisance structure that the analysis claims to remove.

The third level is biological and must match the type of the prediction.  A scalar position score
can be compared with contacts or distances.  A typed edge operator calls for a gauge-aligned
double-mutant matrix.  A background-dependent response calls for the same perturbations across
multiple sequence contexts.  A higher-order residual calls for combinatorial interventions beyond
pairs.  Validation becomes stronger when the endpoint retains the structure of the inferred
object instead of compressing it prematurely.

Falsification is constructive here.  Failure of static agreement can reveal transfer or
misspecification.  Failure of reciprocity can expose directed computation.  Failure of pair
truncation can quantify higher-order dependence.  A theory earns scope not by surviving every
possible test, but by saying precisely what each failure means.
